\textbf{Semantically compatible tasks.} Since topology captures recurrence and movement structure rather than timbre, promising tasks are those where temporal trajectory shape is inherently informative: structural segmentation, recognition of musical form (verse/chorus/bridge), and cover-song identification (timbre changes but structural skeleton remains). A positive result on such tasks, combined with a negative result on genre, would support the proposed explanation separating ``shape'' from ``position.''

\textbf{Cycle as a repetition detector.} A homological representative cycle extracted through Dionysus can be compared with existing segmentation methods based on self-similarity matrices and evaluated as a standalone repetition detector. The cocycle/cycle distinction identified in this work directly informs such research.

\textbf{Embedding interpretability.} Weak topological consistency across spaces (H4) raises the question of how representation topology relates to training objectives and model data. This connects the work to interpretability research on neural audio embeddings.

\textbf{Methodological refinements.} (1)~Systematic selection of delay-embedding parameters, for example via mutual information~\cite{fraser1986independent}, instead of heuristics. (2)~Multivariate and pseudoperiodic surrogates as finer null models. (3)~Higher-dimensional homology ($H_2$, $H_3$), whose musical interpretation remains unexplored. (4)~A stricter H-loop null model using matched random sets by temporal-distance distribution or block/circular permutation.
